\section{第二世纪的边地日常：警报、市场与迁徙}
\label{sec:eastern-han-second-borderlands}

东汉第二世纪的边地问题，不能只由雒阳朝廷发出的“叛”“降”“平定”来命名。
凉州、并州、河西与北方草原相接的郡县，既是军队、牧地与关塞的所在，也是市场、
家户、翻译、向导和地方官吏必须继续生活的地方。警报传到都城后会变成调兵与
赏赐；警报尚未传出之前，居民已在决定是守城、迁徙、交纳粮草、投靠豪右，还是
同邻近部众谈判。史书较容易保存将军和诏书，较少保存这些日常选择，不能让材料
的沉默变成“边民没有行动”的结论。

\subsection{一份军粮如何越过陇山}

羌地冲突反复发生时，雒阳最难解决的未必是发布命令，而是使命令附着在可用的
粮、马、道路和中介上。军粮离开仓廪后，要经过转运、看守、山口、河谷和驻点；
每一处都可能遇到季节、损耗、徭役不足或地方关系改变。将领的奏议常把这些困难
写成应否进兵的理由，不能据此精确计算运输量，却说明“远征”并非一项只在战场
发生的行为。赵充国的屯田、招抚与分化主张之所以值得反复阅读，正在于它承认
持续供给比一次突击更能决定边政能否维持；它也会把土地、劳力与军事保护重新
压在边地社会身上。

东汉的凉州危机继承了这种难题。羌、胡、鲜卑等名称来自不同时刻的军政分类，
不能替代每个部众、城邑或迁徙家庭的共同意图。某些人可能为保住牧地而战，有些
人为了躲避征发而移动，有些人则在汉将、地方豪强与其他首领之间寻找较可接受
的保护。把他们都写作一个与帝国相对的“边族”，反而遮住战争为何能够一次次
重新组织道路、劳役与军队。

\subsection{塞市与关文的两个方向}

边界并非只有战马越过的一条线。河西与北边的塞市、驿传和关文，使商人、使节、
军人和地方居民在同一处处理货物、身份与通行。居延、肩水金关和悬泉置的简牍
主要属于西汉和特定岗位，却能帮助读者理解汉代边防的行政质地：验传、给马、
供食、转递与查核都是国家能力的一部分，也都依赖具体的人按时到位。东汉的材料
保存得不一样，不能把西汉简牍直接当成第二世纪的完整账册；两者相互照亮的是，
没有这些小规模劳动，所谓“守边”便只剩一条朝廷地图。

西域的绿洲关系也需要从这个尺度理解。班超、甘英和都护职名能够标出汉廷曾把
使节与军政联系送到何处，不能画出稳定疆界。绿洲王廷、商队与中间政权有各自
的水源、市场和安全计算；一条路是否开放，常取决于谁愿意提供向导和接待，而
不是都护的名号本身。甘英未至大秦留下的并非“失败的外交逸事”，而是远距离
信息经过安息、海路和商人中介后仍无法由皇帝命令校正的证据。

\subsection{迁徙留下的郡县}

军费和战争的后果最终会回到人户。有人被编入新的屯戍，有人因缺粮离开原有土地，
有人在新的城邑、宗族或主从关系中寻求保护。地方官若要恢复赋役和治安，必须先
决定谁可被登记、谁有资格申诉、谁应承担转运；这些决定又会影响人们是否愿意
回到旧籍。正史中的“流民”“降人”或“徙民”不能抹去这种差异，它们只是国家
在最需要分类时给出的名称。

到一八四年以后，州郡和边将能自行筹兵的原因正在这里累积。黄巾与凉州战事并
非同一事件，却都迫使地方在常规拨付无法及时抵达时寻找粮饷、武器和可用人员。
州牧扩权不是割据的预言，而是面对多处压力时的一次应急选择；它的后果则是仓储、
道路和人与人之间的保护关系越来越能够服务区域中心。董卓等人的军政资源进入
宫廷危机，正是这条长期边地压力线在都城取得入口的结果。

\begin{turningpoint}[边防为何改变了地方政治]
边地的军费暂时解决的是城塞、道路和使节的安全；它要求地方持续提供粮、马、
劳力和信息。中枢若不能以稳定规则支付这些成本，能筹集资源的人便取得新的谈判
位置。第二世纪的迁徙、塞市与军镇因此不是汉末政治的背景，而是州郡军事化得以
发生的社会条件。
\end{turningpoint}

\subsection{陇右的春耕为何也在战事之中}

西羌在建初元年于陇西、金城一带的反叛，并不是一张已经写好敌我双方的战场图。
\eventref{EH-0076-QIANG-REBELLION}{这场边地冲突}的年份和它所引出的长期压力可以
核定，至于哪些部众在何时以何种关系卷入、各县居民如何站队，则不能由“羌”
这一朝廷分类一笔写完。陇右的山谷、渡口和牧地既是军队行进的空间，也是耕作、
放牧、交换和婚姻关系延续的地方。一处警报传来，地方家庭面对的未必是“支持”
或“反对”朝廷这样的选择，而可能是今年是否还能下种、牲畜能否保住、亲属能否
经过山口、县城能否替他们裁断争执。

春耕使边战尤其难以从地方生活中切开。役夫和畜力被抽去转运，留在村落的人就
要重新安排耕作；道路被军队优先使用，市场上的谷物、盐和布帛也更难按原先节奏
流动。军事行动若在一处取得战果，未必就能使下一季的播种恢复；反之，一处村落
能够继续耕作，也不表示它已经脱离威胁。正史往往以平叛、受降和奏捷划分时间，
而地方的时间还受雨水、草场、种子和欠下的劳役支配。把两种时间并置，才能避免
让边民只在将领出场时才有历史。

朝廷尝试以迁徙、招抚、屯戍和更严密的守备来处理这种局面。这些措施各有不同
的即时目的：迁人可为某处补劳力或减少眼前冲突，招抚可拆开敌对关系，屯戍可让
军队靠近粮源，守备则试图保住通道。没有一项是无代价的。迁出者失去原有田地和
亲属网络，接纳地要重新分配水源与役使；被招抚者未必从此脱离旧有关系；屯戍
所需的土地、种子和保护又会反过来加重地方负担。我们不知道这些方案在每一座
城、每一道山谷的实际效果，正因如此，更不应把朝廷的制度名称误读为已完成的
社会事实。

\subsection{一次“平定”以后，谁来使用留下的道路}

\eventanchor{EH-0167-DUAN-JIONG-QIANG}{段颎在陇右连续进军}使朝廷短期内恢复了部分
通道和郡县的控制。战役序列大体可知，传世叙事中的斩获、迁徙规模和“尽平”的
语言却不能当作无误统计。尤其需要警惕的是，把军事结果写成边地社会已经重新
安定。道路能够重开，首先意味着军队、税粮、使节与商旅又有一条可走的路；它
不能说明被击散、降附或迁徙的人已在同一制度内取得安全，也不能保证沿线家庭
不再承担征发和报复的风险。

严酷围剿留下的政治问题，常在战报之后才开始显形。郡县需要重新登记人户，
将领需要安置士卒和降附者，地方有力者需要决定是否接纳来者、如何处理无主的
田地和牲畜。对朝廷来说，恢复关隘和税源的诱惑很大；对地方社会来说，新的
登记和守备可能意味着旧有的依附、债务和土地争执被重新打开。所谓“恢复秩序”
并不是把一个静止的旧世界放回原位，而是让战后已经改变的人、物和关系进入新的
可管理形式。

这也是为什么边地的军政人物会积累超过一次战役的影响。他们熟悉粮道、马匹、
降附者和地方首领，能够在中枢不在场时判断何处可守、何处可借粮、何处必须让步。
这种知识能服务于国家，也会成为个人和区域网络的资源。不能据此把每位边将写成
预谋割据者；他们取得的行动余地，首先来自国家需要有人把遥远的命令转成眼前的
安排。当同样的压力在下一轮动乱中出现，旧的知识、宾客和供给关系便不会从零
开始。

\subsection{北边的马、市与一场难以结算的出师}

\eventref{EH-0177-XIANBEI-RAIDS}{熹平六年对鲜卑的出击受挫}，把北边问题再一次
带回朝廷。侵扰与出师的事实可以确认，但不能由此画出檀石槐领导下所有部众的
整齐组织，也不能以传世数字衡量双方损失。北边的难处不只在于敌军能否被击退。
马匹、草场、塞市、降附者、守兵家属和通往内地的道路，都是随战事一起被重新
安排的资源。朝廷若只以一次远征解决问题，便要承担长距离补给的成本；若减少
出击，又须考虑边郡居民、互市关系和邻近集团如何理解这一变化。

市场在这里不是战争的反面。互市和私下往来能够为边地提供日用品、牲畜和信息，
同时也使官府担心军用物资、人员和消息越过关口。严加盘查会提高流动的成本，
完全放开又难以维持守备的界线。县吏、关吏、商人、牧民和使者必须在这种不确定
中反复辨验身份和货物。我们无法把每一笔交易写出来，却可以看见：边界之所以
能被维持，靠的并不只是烽火和兵器，也靠许多关于通行、交割与担保的日常决定。

到灵帝末年，北边、凉州和中原的危机彼此牵动，不能把其中任一处当成另一处的
单纯背景。远征失败会改变朝廷对军费与将领的判断，地方征发会改变家庭对迁徙和
保护的判断，边市和道路的紧张则会影响区域中心取得马匹与消息的能力。第二世纪
的边地因而不是汉末突然崩解前静止的边缘；它是一处持续学习如何在供给不足、
关系多变和分类不稳中维持生活的地方。后来州郡军事化所利用的，正是这种长期
积累的道路、人员与不得不就地决断的能力。
